\subsection*{29. Real Price Tests and the Risk of Confounding}

Suppose 2 comparable groups receive the same specified offer at ₹24,000 and ₹36,000. Track eligible impressions, qualified inquiries, applications, actual payment, refunds, and later outcomes. If groups are genuinely comparable, the test may reveal a causal difference over the observed price range. But differences in channel, instructor identity, season, urgency, or marketing narrative can invalidate the comparison.

Publicly inconsistent pricing can also cause fairness problems. When randomized price differences for identical customers would be inappropriate, test transparent differences in package scope, different launch cohorts with clearly stated policies, or targeted sponsorships. Record the limitations of sequential comparisons: willingness to pay can shift as reputation and demand change.

One fundamental rule: \textbf{never accept a paid pre-order for an offering you cannot credibly fulfill under the stated terms}. Experiments are not exemptions from consumer trust.

\subsection*{30. Use Refusals as Diagnostic Evidence}

A refusal usually includes more information than ``price too high.'' Categorize it across 7 diagnostic categories:

\paragraph*{1. No perceived incremental value:} suggests the wrong job or poor differentiation.
\paragraph*{2. Cannot verify quality:} suggests a proof or credence problem.
\paragraph*{3. Cannot afford safely:} suggests payer mismatch or a different service architecture.
\paragraph*{4. No time:} suggests nonmonetary adoption sacrifice is too steep.
\paragraph*{5. Approval required:} suggests an unaddressed B2B buying-process obstacle.
\paragraph*{6. Not urgent:} suggests poor timing or lack of an immediate trigger.
\paragraph*{7. Adequate free alternative:} suggests the paid feature fails to deliver meaningful additional progress.

Ask what alternative the buyer will now use. Request a concrete example of what evidence would change the decision. Resist offering an immediate discount: the 1st discount can conceal the very objection the founder needed to learn about. It can also anchor future customers to a lower fee before the value proposition has been improved.

\subsection*{31. How Much Evidence Is Enough?}

No fixed sample size establishes a market's willingness-to-pay distribution. Statistical precision depends on effect size, conversion rate, selection, noise, and the decision at stake. A founder may only need evidence that a dozen qualified buyers willingly purchase an offer at a viable price to justify a 2nd bounded pilot. That same evidence cannot establish the optimum price for a national market.

Instead of universal thresholds, precommit to decisions. For example: ``We will repeat at this fee if at least 10 well-qualified independent buyers pay without exceptional pressure, the defined delivery outcome is achieved, refunds remain manageable, and normalized contribution per cohort is positive.'' These numbers are \textbf{illustrative decision criteria}, not industry benchmarks. A failing threshold could call for changing the customer, price, proof, or product: not automatically abandoning the institution.

The best research process does not eliminate uncertainty. It purchases sufficient reduction in uncertainty to justify the next commitment.
